\documentclass[10pt,a4paper]{amsart}
\usepackage[margin=19mm]{geometry}
\usepackage{amsmath,amssymb,mathtools}
\usepackage[scr=rsfs,cal=euler]{mathalfa}
\usepackage{xfrac}
\usepackage[unicode,hidelinks,bookmarks=false]{hyperref}
\newcommand{\ie}{i.e.\ }
\newcommand{\cf}{cf.\ }
\newcommand{\resp}{respectively\ }
\newcommand{\Subsection}{\subsection}
\providecommand{\nobreakdash}{\nobreak\hbox{-}}
\setlength{\emergencystretch}{2em}
\multlinegap0pt
\usepackage{amssymb}
\newtheorem{thm}{Theorem}
\newcommand{\D}{\mathbb{D}}
\newcommand{\bsa}{{\boldsymbol a}}
\newcommand{\bsb}{{\boldsymbol b}}
\newcommand{\bsx}{{\boldsymbol x}}
\newcommand{\bx}{{\rm box}}
\newcommand{\cA}{\mathcal{A}}
\newcommand{\cP}{\mathcal{P}}
\newcommand{\ex}{{\rm ext}}
\newcommand{\rd}{\,{\rm d}}
\newcommand{\sgn}{{\rm sign}}
\title{Extreme $L_p$ discrepancy, numerical integration and the curse of dimensionality}
\author{Erich Novak, Friedrich Pillichshammer}
\date{Source: arXiv:2602.19760v3 (CC BY 4.0)}
\begin{document}
\maketitle

\noindent\emph{Source and licence.} Extreme $L_p$ discrepancy, numerical integration and the curse of dimensionality. Version arXiv:2602.19760v3 (CC BY 4.0), distributed by arXiv under the Creative Commons Attribution 4.0 International licence, http://creativecommons.org/licenses/by/4.0/, as recorded in the arXiv licence metadata for this exact version and shown on the abstract page https://arxiv.org/abs/2602.19760v3. Full licence notice: \texttt{licenses/arxiv-2602.19760-CC-BY-4.0.txt}.\par\smallskip

\noindent\textit{Editorial scope.} Counted unit: the article's thm thm:main (source0.tex lines 247--253). The article's complete original proof of it is delivered with it (source0.tex lines 256--324, one \texttt{proof} environment of 69 lines). No mathematics is written, repaired or re-derived: the statement and the proof are the article's own lines, in the article's own notation, and no retained line is substituted or re-flowed.  Retained uncounted as the source-local context the proof needs: lines 186--189; lines 191--193; lines 238--241.  Nothing was omitted from any retained span, so the package carries no omission record.  No retained line contains a citation, so the wrapper carries no bibliography and no bibliography entry.  No retained line contains a figure environment, an image-inclusion command or drawing code, so no image is delivered and the package declares no asset dependency.  Editorial formatting only, in addition: an AMS \texttt{amsart} preamble that restates the article's own declarations and macros used by the retained lines, namely the environments \texttt{thm} on separate counters, together with the standard packages those lines need.  The retained environments print in this document as Theorem 1 in order of appearance, whereas the article prints the same items with the numbering of its own full document.  The complete native source of the article as distributed in the arXiv e-print for this exact version is bundled unchanged as \texttt{sources/195139/source0.tex}.
\begin{equation}\label{eq:rep}
(T_d c)(\bsx):=\int_{\D_d} c(\bsa,\bsb)\, {\bf 1}_{[\bsa,\bsb]}(\bsx) \rd (\bsa,\bsb)
\quad\text{for }c\in L_q(\D_d).
\end{equation}

\begin{equation}\label{eq:norm-rep}
\|f\|_{F_{d,q},\bx}:=\inf\left\{\|c\|_{L_q(\D_d)}:\;c\in L_q(\D_d), \ f=T_dc\right\}.
\end{equation}

\begin{equation}\label{eq:L-mu}
L(f)=\int_{[0,1]^d} f(\bsx) \rd\left(\lambda-\sum_{k=1}^N c_k\delta_{\bsx_k}\right)(\bsx)
= -\int_{[0,1]^d} f(\bsx) \rd\mu(\bsx).
\end{equation}

\begin{thm}[Discrepancy-integration duality]\label{thm:main}
Let $p,q \in [1,\infty ]$ be H\"older conjugates. For every point set 
$\cP=\{\bsx_1,\ldots,\bsx_N\}$ in $[0,1)^d$ and arbitrary corresponding real weights $\cA=\{c_1,\ldots,c_N\}$ we have
\[
e(Q_{\cP,\cA};F_{d,q}) \;=\; L_{p,N}^{\ex}(\cP,\cA).
\]
\end{thm}

\begin{proof}
Let $f\in F_{d,q}$ and choose $c\in L_q(\D_d)$ representing $f$ as in \eqref{eq:rep}.
Using \eqref{eq:L-mu} and Fubini, we compute
\begin{align*}
L(f)
&= -\int_{[0,1]^d}  f(\bsx) \rd\mu(\bsx)
= -\int_{[0,1]^d} \left(\int_{\D_d} c(\bsa,\bsb) \,{\bf 1}_{[\bsa,\bsb]}(\bsx)\rd (\bsa,\bsb)\right) \rd \mu(\bsx)\\
&= -\int_{\D_d} c(\bsa,\bsb)\left(\int_{[0,1]^d} {\bf 1}_{[\bsa,\bsb]}(\bsx)\,d\mu(\bsx)\right) \rd (\bsa,\bsb)\\
&= -\int_{\D_d} c(\bsa,\bsb) \,\mu([\bsa,\bsb]) \rd (\bsa,\bsb)\\
&= -\int_{\D_d} c(\bsa,\bsb)\,\Delta_{\cP,\cA}([\bsa,\bsb]) \rd (\bsa,\bsb).
\end{align*}
Now, by Hölder's inequality,
\[
|L(f)|
\le \|c\|_{L_q(\D_d)}\,\|\Delta_{\cP,\cA}\|_{L_p(\D_d)}
= \|c\|_{L_q(\D_d)}\,L_{p,N}^{\ex}(\cP,\cA).
\]
Taking the infimum over all representations of $f$ yields
\[
|L(f)| \le \|f\|_{F_{d,q},\bx}\,L_{p,N}^{\ex}(\cP,\cA).
\]
Hence $$e(Q_{\cP,\cA};F_{d,q})\le L_{p,N}^{\ex}(\cP,\cA).$$

It remains to prove equality. First we assume $p,q \in (1, \infty)$ and discuss the needed modifications 
for the endpoints later. If $\Delta_{\cP,\cA} \equiv 0$ then $L_{p,N}^{\ex}(\cP,\cA)=0$ and also $L(f)\equiv 0$, 
so equality holds. Assume now $\Delta_{\cP,\cA}\not\equiv 0$. Define 
\begin{equation}\label{eq:cstar}
c^{\ast}(\bsa,\bsb):=\frac{|\Delta_{\cP,\cA}([\bsa,\bsb])|^{p-2} \Delta_{\cP,\cA}([\bsa,\bsb])}{\|\Delta_{\cP,\cA}\|_{L_p(\D_d)}^{p-1}}
\quad\text{so that}\quad
\|c^{\ast}\|_{L_q(\D_d)}=1.
\end{equation}
Let $f^\ast=T_d c^\ast$ with coefficient $c^\ast$. Then $\|f^\ast\|_{F_{d,q},\bx}\le 1$ by definition of the norm \eqref{eq:norm-rep}, and
\begin{align*}
L(f^\ast)
= & -\int_{\D_d} c^\ast(\bsa,\bsb)\,\Delta_{\cP,\cA}([\bsa,\bsb])\rd (\bsa,\bsb) \\
= & -\frac{1}{\|\Delta_{\cP,\cA}\|_{L_p(\D_d)}^{p-1}}\int_{\D_d}|\Delta_{\cP,\cA}(\bsa,\bsb)|^p \rd (\bsa,\bsb)\\
= & -\|\Delta_{\cP,\cA}\|_{L_p(\D_d)}.
\end{align*}
Therefore $|L(f^\ast)|=\|\Delta_{\cP,\cA}\|_{L_p(\D_d)}=L_{p,N}^{\ex}(\cP,\cA)$ and hence
$$e(Q_{\cP,\cA};F_{d,q})\ge L_{p,N}^{\ex}(\cP,\cA).$$ Combining with the upper bound yields the identity.

For the case $p=1$ we take $c^*(\bsa,\bsb) = \sgn (\Delta_{\cP,\cA} (\bsa,\bsb))$ and obtain 
again $|L(f^\ast)|=\|\Delta_{\cP,\cA}\|_{L_p(\D_d)}=L_{p,N}^{\ex}(\cP,\cA)$.
Assume now that  $p=\infty$ and  $q=1$. Let
\[
M:=\Vert \Delta_{\cP,\cA} \Vert_{L_{\infty}(\D_d)}.
\]
For any $\varepsilon>0$, define
\[
E_\varepsilon
:=\{(\bsa,\bsb): |\Delta_{\cP,\cA}(\bsa,\bsb)| \ge M-\varepsilon\}.
\]
Assume $m(E_\varepsilon)>0$.
Define
\[
c_\varepsilon (\bsa, \bsb) 
=\frac{ \sgn (\Delta_{\cP,\cA} (\bsa, \bsb) )\,  {\mathbf 1}_{E_\varepsilon} (\bsa, \bsb) }{m(E_\varepsilon)}.
\]
Then $\Vert c_\varepsilon\Vert_{L_1}=1$ and
\[
\int_{\D_d}  c_\varepsilon (\bsa,\bsb)\,\Delta_{\cP,\cA}([\bsa,\bsb])\rd (\bsa,\bsb )
\ge M-\varepsilon.
\]
Letting $\varepsilon\to0$ gives
\[
e(Q_{\cP,\cA};F_{d,1})
=\Vert  \Delta_{\cP,\cA}\Vert_{L_\infty}=L_{\infty,N}^{\ex}(\cP,\cA).
\]
\end{proof}

\end{document}
